%==============================================================================
%  ENGR 2405 - Electrical Circuits I
%  Lab 9 - Sinusoids & Phasors
%
%  Compile with pdfLaTeX (or XeLaTeX / LuaLaTeX).
%==============================================================================
\documentclass[11pt]{article}

%------------------------- Packages -------------------------------------------
\usepackage[utf8]{inputenc}
\usepackage[T1]{fontenc}
\usepackage[margin=1in]{geometry}     % page margins
\usepackage{amsmath,amssymb}          % math
\usepackage{graphicx}                 % \includegraphics (kept for optional schematics)
\usepackage{float}                    % [H] figure placement
\usepackage{booktabs}                 % nicer tables
\usepackage{siunitx}                  % \SI{5}{\ohm}, \SI{10}{\milli\second}, etc.
\usepackage{listings}                 % SciLab code blocks
\usepackage{xcolor}
\usepackage{caption}
\usepackage[hidelinks]{hyperref}
\usepackage{titlesec}

%------------------------- Section styling ------------------------------------
\titleformat{\section}{\Large\bfseries}{}{0pt}{}
\titlespacing*{\section}{0pt}{1.4em}{0.6em}

%------------------------- Code listing style ---------------------------------
\definecolor{codebg}{RGB}{245,245,245}
\definecolor{codecomment}{RGB}{0,128,0}
\definecolor{codekw}{RGB}{0,0,180}
\lstset{
    backgroundcolor=\color{codebg},
    basicstyle=\ttfamily\small,
    breaklines=true,
    frame=single,
    framesep=4pt,
    keywordstyle=\color{codekw}\bfseries,
    commentstyle=\color{codecomment}\itshape,
    showstringspaces=false,
    captionpos=b,
    tabsize=2,
}
\definecolor{codestr}{RGB}{163,21,21}      % strings
\definecolor{codenum}{RGB}{120,120,120}    % numbers

% --- SciLab (custom: listings has no built-in SciLab language) -------------
\lstdefinelanguage{Scilab}{
    morekeywords={
        function,endfunction,if,then,else,elseif,end,for,while,do,
        select,case,break,continue,return,clc,clear,disp,mprintf,
        printf,zeros,ones,eye,linspace,inv,det,sum,length,size,abs,
        sqrt,real,imag,exp,sin,cos,tan,plot,plot2d,xlabel,ylabel,
        title,legend,grid,poly,roots,find,pause,error,warning
    },
    sensitive=true,
    morecomment=[l]{//},          % line comment
    morecomment=[s]{/*}{*/},      % block comment
    morestring=[b]",              % double-quoted strings
    morestring=[b]',              % single-quoted strings
}
\lstdefinestyle{scilab}{
    language=Scilab,
    stringstyle=\color{codestr},
}

%==============================================================================
\begin{document}

%------------------------- Title block ----------------------------------------
\begin{center}
    {\LARGE\bfseries Lab 9 --- Sinusoids \& Phasors}\\[0.8em]
    {\large Abdon Morales}\\[0.3em]
    {\large \today}
\end{center}
\vspace{1em}

%------------------------------------------------------------------------------
\section*{Description}

This lab looks at how series RC, RL, and RLC circuits behave when they are driven
by a sine wave instead of a DC source. Each circuit is built on the trainer,
driven with a \SI{10}{\volt} peak-to-peak sine source, and the frequency is swept
from \SI{100}{\hertz} to \SI{10}{\kilo\hertz}. At each frequency we measure the
amplitude of the reactive element voltage and its phase difference relative to the
source.

The point is to see impedance behave the way resistance does not. Impedance is a
complex number $Z = R + jX$, so the series combination rules from DC still apply,
but the reactance $X$ depends on frequency: $Z_C = 1/(j\omega C)$ shrinks as
frequency goes up, while $Z_L = j\omega L$ grows. Because the impedances are
complex, the element voltages come out shifted in phase from the source, and the
voltage divider that sets each element voltage changes as the frequency changes.
We express all of the results as phasors and compare the measured magnitude and
phase against the calculated values.

%------------------------------------------------------------------------------
\section*{Procedure}

\begin{enumerate}
    \item Built a series RC circuit with the \SI{4.7}{\kilo\ohm} resistor and the
          \SI{0.1}{\micro\farad} capacitor across the sine wave source on the
          XK700 trainer.
    \item Set the source amplitude to \SI{10}{\volt} peak-to-peak and verified it
          on the oscilloscope. Kept channel 1 on the source and channel 2 on the
          element being measured.
    \item Stepped the source frequency through \SI{100}{\hertz},
          \SI{500}{\hertz}, \SI{1}{\kilo\hertz}, \SI{5}{\kilo\hertz}, and
          \SI{10}{\kilo\hertz}. At each step we rechecked the source amplitude
          (the generator output sags a little when the load impedance drops) and
          recorded $V_C$ peak-to-peak and the phase difference $\phi_C$ between
          $V_C$ and the source.
    \item Measured phase from the time shift between the two zero crossings,
          $\phi = -360^\circ \, \Delta t \, f$, taking the sign from whether the
          element voltage crossed before or after the source.
    \item Rebuilt the circuit as a series RL using the \SI{100}{\ohm} resistor and
          the \SI{33}{\milli\henry} inductor, and repeated steps 2--4 recording
          $V_L$ and $\phi_L$.
    \item Rebuilt the circuit as a series RLC using the \SI{1}{\kilo\ohm}
          resistor, the \SI{0.33}{\micro\farad} capacitor, and the
          \SI{10}{\milli\henry} inductor. Repeated steps 2--4 at each frequency,
          this time recording both $V_C$, $\phi_C$ and $V_L$, $\phi_L$. Moving the
          scope probe between the capacitor and the inductor required
          re-referencing ground, so we reordered the series elements as needed to
          keep the element under test next to the grounded side.
\end{enumerate}

%------------------------------------------------------------------------------
\section*{Calculations}

All three circuits are a single series loop, so the total impedance is just the
sum of the element impedances and each element voltage follows from the impedance
divider. With the source taken as the phase reference,
$\mathbf{V}_s = 10\angle 0^\circ$ V (peak-to-peak),
\[
    \mathbf{Z}_C = \frac{1}{j\omega C} = \frac{-j}{\omega C},
    \qquad
    \mathbf{Z}_L = j\omega L,
    \qquad
    \omega = 2\pi f .
\]

\noindent\textbf{RC circuit} ($R = \SI{4.7}{\kilo\ohm}$, $C = \SI{0.1}{\micro\farad}$):
\[
    \mathbf{Z}_T = R + \frac{1}{j\omega C},
    \qquad
    \mathbf{V}_C = \mathbf{V}_s \frac{\mathbf{Z}_C}{\mathbf{Z}_T},
    \qquad
    \mathbf{V}_R = \mathbf{V}_s \frac{R}{\mathbf{Z}_T}.
\]

\noindent\textbf{RL circuit} ($R = \SI{100}{\ohm}$, $L = \SI{33}{\milli\henry}$):
\[
    \mathbf{Z}_T = R + j\omega L,
    \qquad
    \mathbf{V}_L = \mathbf{V}_s \frac{\mathbf{Z}_L}{\mathbf{Z}_T}.
\]

\noindent\textbf{RLC circuit} ($R = \SI{1}{\kilo\ohm}$, $C = \SI{0.33}{\micro\farad}$,
$L = \SI{10}{\milli\henry}$):
\[
    \mathbf{Z}_T = R + j\!\left(\omega L - \frac{1}{\omega C}\right),
    \qquad
    \mathbf{V}_C = \mathbf{V}_s \frac{\mathbf{Z}_C}{\mathbf{Z}_T},
    \qquad
    \mathbf{V}_L = \mathbf{V}_s \frac{\mathbf{Z}_L}{\mathbf{Z}_T}.
\]

The series RLC has a resonant frequency where the two reactances cancel:
\[
    f_0 = \frac{1}{2\pi\sqrt{LC}}
        = \frac{1}{2\pi\sqrt{(\SI{10}{\milli\henry})(\SI{0.33}{\micro\farad})}}
        = \SI{2.77}{\kilo\hertz},
\]
which falls between the \SI{1}{\kilo\hertz} and \SI{5}{\kilo\hertz} test points.

% --- SciLab program ---
\begin{lstlisting}[style=scilab, caption={SciLab program: theoretical impedances and phasor voltages}]
// Lab 9 - Sinusoids & Phasors
// Series RC, RL, RLC: total impedance and element voltage phasors

f  = [100 500 1000 5000 10000];   // Hz
Vs = 10;                          // Vpp, reference phase 0
w  = 2*%pi*f;

function polarprint(name, z)
    mprintf("  %-4s = %9.3f  angle %8.2f deg\n", name, abs(z), atan(imag(z),real(z))*180/%pi);
endfunction

// ---------------- RC: R = 4.7k, C = 0.1u ----------------
R = 4700; C = 0.1e-6;
mprintf("\n=== RC circuit: R = 4.7k, C = 0.1u ===\n");
for k = 1:length(f)
    Zc = 1/(%i*w(k)*C);
    Zt = R + Zc;
    Vc = Vs*Zc/Zt;
    Vr = Vs*R/Zt;
    mprintf("f = %6.0f Hz\n", f(k));
    polarprint("Zc", Zc); polarprint("Zt", Zt);
    polarprint("Vc", Vc); polarprint("Vr", Vr);
end

// ---------------- RL: R = 100, L = 33m ----------------
R = 100; L = 33e-3;
mprintf("\n=== RL circuit: R = 100, L = 33m ===\n");
for k = 1:length(f)
    Zl = %i*w(k)*L;
    Zt = R + Zl;
    Vl = Vs*Zl/Zt;
    Vr = Vs*R/Zt;
    mprintf("f = %6.0f Hz\n", f(k));
    polarprint("Zl", Zl); polarprint("Zt", Zt);
    polarprint("Vl", Vl); polarprint("Vr", Vr);
end

// ---------------- RLC: R = 1k, C = 0.33u, L = 10m ----------------
R = 1000; C = 0.33e-6; L = 10e-3;
f0 = 1/(2*%pi*sqrt(L*C));
mprintf("\n=== RLC circuit: R = 1k, C = 0.33u, L = 10m ===\n");
mprintf("resonant frequency f0 = %.1f Hz\n", f0);
for k = 1:length(f)
    Zc = 1/(%i*w(k)*C);
    Zl = %i*w(k)*L;
    Zt = R + Zc + Zl;
    Vc = Vs*Zc/Zt;
    Vl = Vs*Zl/Zt;
    mprintf("f = %6.0f Hz\n", f(k));
    polarprint("Zl", Zl); polarprint("Zc", Zc); polarprint("Zt", Zt);
    polarprint("Vc", Vc); polarprint("Vl", Vl);
end
\end{lstlisting}

% --- SciLab runtime printout ---
\begin{lstlisting}[caption={SciLab runtime output}]
=== RC circuit: R = 4.7k, C = 0.1u ===
f =    100 Hz
  Zc   = 15915.494  angle   -90.00 deg
  Zt   = 16594.968  angle   -73.55 deg
  Vc   =     9.591  angle   -16.45 deg
  Vr   =     2.832  angle    73.55 deg
f =    500 Hz
  Zc   =  3183.099  angle   -90.00 deg
  Zt   =  5676.453  angle   -34.11 deg
  Vc   =     5.608  angle   -55.89 deg
  Vr   =     8.280  angle    34.11 deg
f =   1000 Hz
  Zc   =  1591.549  angle   -90.00 deg
  Zt   =  4962.160  angle   -18.71 deg
  Vc   =     3.207  angle   -71.29 deg
  Vr   =     9.472  angle    18.71 deg
f =   5000 Hz
  Zc   =   318.310  angle   -90.00 deg
  Zt   =  4710.767  angle    -3.87 deg
  Vc   =     0.676  angle   -86.13 deg
  Vr   =     9.977  angle     3.87 deg
f =  10000 Hz
  Zc   =   159.155  angle   -90.00 deg
  Zt   =  4702.694  angle    -1.94 deg
  Vc   =     0.338  angle   -88.06 deg
  Vr   =     9.994  angle     1.94 deg

=== RL circuit: R = 100, L = 33m ===
f =    100 Hz
  Zl   =    20.735  angle    90.00 deg
  Zt   =   102.127  angle    11.71 deg
  Vl   =     2.030  angle    78.29 deg
  Vr   =     9.792  angle   -11.71 deg
f =    500 Hz
  Zl   =   103.673  angle    90.00 deg
  Zt   =   144.042  angle    46.03 deg
  Vl   =     7.197  angle    43.97 deg
  Vr   =     6.942  angle   -46.03 deg
f =   1000 Hz
  Zl   =   207.345  angle    90.00 deg
  Zt   =   230.200  angle    64.25 deg
  Vl   =     9.007  angle    25.75 deg
  Vr   =     4.344  angle   -64.25 deg
f =   5000 Hz
  Zl   =  1036.726  angle    90.00 deg
  Zt   =  1041.537  angle    84.49 deg
  Vl   =     9.954  angle     5.51 deg
  Vr   =     0.960  angle   -84.49 deg
f =  10000 Hz
  Zl   =  2073.451  angle    90.00 deg
  Zt   =  2075.861  angle    87.24 deg
  Vl   =     9.988  angle     2.76 deg
  Vr   =     0.482  angle   -87.24 deg

=== RLC circuit: R = 1k, C = 0.33u, L = 10m ===
resonant frequency f0 = 2770.5 Hz
f =    100 Hz
  Zl   =     6.283  angle    90.00 deg
  Zc   =  4822.877  angle   -90.00 deg
  Zt   =  4919.307  angle   -78.27 deg
  Vc   =     9.804  angle   -11.73 deg
  Vl   =     0.013  angle   168.27 deg
f =    500 Hz
  Zl   =    31.416  angle    90.00 deg
  Zc   =   964.575  angle   -90.00 deg
  Zt   =  1367.767  angle   -43.02 deg
  Vc   =     7.052  angle   -46.98 deg
  Vl   =     0.230  angle   133.02 deg
f =   1000 Hz
  Zl   =    62.832  angle    90.00 deg
  Zc   =   482.288  angle   -90.00 deg
  Zt   =  1084.409  angle   -22.76 deg
  Vc   =     4.447  angle   -67.24 deg
  Vl   =     0.579  angle   112.76 deg
f =   5000 Hz
  Zl   =   314.159  angle    90.00 deg
  Zc   =    96.458  angle   -90.00 deg
  Zt   =  1023.423  angle    12.28 deg
  Vc   =     0.942  angle  -102.28 deg
  Vl   =     3.070  angle    77.72 deg
f =  10000 Hz
  Zl   =   628.319  angle    90.00 deg
  Zc   =    48.229  angle   -90.00 deg
  Zt   =  1156.073  angle    30.12 deg
  Vc   =     0.417  angle  -120.12 deg
  Vl   =     5.435  angle    59.88 deg
\end{lstlisting}

\subsection*{Theoretical Results}

\begin{table}[H]
    \centering
    \caption{RC circuit, theoretical ($R = \SI{4.7}{\kilo\ohm}$, $C = \SI{0.1}{\micro\farad}$)}
    \begin{tabular}{rcccc}
        \toprule
        $f$ (Hz) & $|Z_T|$ (\si{\ohm}) & $\angle Z_T$ (\si{\degree})
                 & $V_C$ (Vpp) & $\phi_C$ (\si{\degree}) \\
        \midrule
        100    & 16595 & $-73.55$ & 9.591 & $-16.45$ \\
        500    &  5676 & $-34.11$ & 5.608 & $-55.89$ \\
        1000   &  4962 & $-18.71$ & 3.207 & $-71.29$ \\
        5000   &  4711 & $-3.87$  & 0.676 & $-86.13$ \\
        10000  &  4703 & $-1.94$  & 0.338 & $-88.06$ \\
        \bottomrule
    \end{tabular}
\end{table}

\begin{table}[H]
    \centering
    \caption{RL circuit, theoretical ($R = \SI{100}{\ohm}$, $L = \SI{33}{\milli\henry}$)}
    \begin{tabular}{rcccc}
        \toprule
        $f$ (Hz) & $|Z_T|$ (\si{\ohm}) & $\angle Z_T$ (\si{\degree})
                 & $V_L$ (Vpp) & $\phi_L$ (\si{\degree}) \\
        \midrule
        100    &  102.1 & $11.71$ & 2.030 & $78.29$ \\
        500    &  144.0 & $46.03$ & 7.197 & $43.97$ \\
        1000   &  230.2 & $64.25$ & 9.007 & $25.75$ \\
        5000   & 1041.5 & $84.49$ & 9.954 & $5.51$  \\
        10000  & 2075.9 & $87.24$ & 9.988 & $2.76$  \\
        \bottomrule
    \end{tabular}
\end{table}

\begin{table}[H]
    \centering
    \caption{RLC circuit, theoretical ($R = \SI{1}{\kilo\ohm}$, $C = \SI{0.33}{\micro\farad}$, $L = \SI{10}{\milli\henry}$)}
    \begin{tabular}{rcccccc}
        \toprule
        $f$ (Hz) & $|Z_T|$ (\si{\ohm}) & $\angle Z_T$ (\si{\degree})
                 & $V_C$ (Vpp) & $\phi_C$ (\si{\degree})
                 & $V_L$ (Vpp) & $\phi_L$ (\si{\degree}) \\
        \midrule
        100    & 4919 & $-78.27$ & 9.804 & $-11.73$  & 0.013 & $168.27$ \\
        500    & 1368 & $-43.02$ & 7.052 & $-46.98$  & 0.230 & $133.02$ \\
        1000   & 1084 & $-22.76$ & 4.447 & $-67.24$  & 0.579 & $112.76$ \\
        5000   & 1023 & $12.28$  & 0.942 & $-102.28$ & 3.070 & $77.72$  \\
        10000  & 1156 & $30.12$  & 0.417 & $-120.12$ & 5.435 & $59.88$  \\
        \bottomrule
    \end{tabular}
\end{table}

%------------------------------------------------------------------------------
\section*{Measurements}

Measurements were taken on a Global Specialties PB-503A trainer (the course
handout specifies an XK700) using a Tektronix TDS 1002 oscilloscope. The TDS 1002
has no automatic phase measurement, so every phase angle was obtained from time
cursors: the period $T$ was read between successive rising zero crossings of the
source on channel 1, the shift $\Delta t$ was read between the rising zero
crossing of the source and the nearest rising zero crossing of the element
voltage on channel 2, and the angle was computed as
\[
    \phi = \pm\, 360^\circ \, \frac{\Delta t}{T}.
\]
Taking $T$ from the trace rather than from the generator dial removes the
generator's frequency accuracy from the result. The sign was assigned from which
trace crossed first; a negative angle means the element voltage lags the source.
Both channels were DC coupled and their zero references were aligned to the
centre graticule line before each set of readings.

Two departures from the ideal model were identified during the lab and are
carried through the tables below. First, the PB-503A function generator has an
output impedance of roughly \SI{600}{\ohm} rather than the \SI{50}{\ohm} assumed
for a bench generator, so it could not hold \SI{10}{\volt} peak-to-peak into the
low-impedance RL circuit. Fitting the measured source amplitude at node A against
a divider model over the five RL test frequencies gives
$R_g \approx \SI{575}{\ohm}$ with an rms residual of \SI{0.07}{\volt}. Because
channel 1 was probed at the circuit node and not at the generator terminals,
$R_g$ falls outside the measured loop and affects only the absolute amplitudes,
not the voltage ratios or the phase angles. Second, the \SI{33}{\milli\henry}
inductor carries a winding resistance of approximately \SI{43}{\ohm}, which is
comparable to $|Z_L| = \SI{20.7}{\ohm}$ at \SI{100}{\hertz}. A DCR column is
therefore included alongside the ideal calculation wherever it changes the
prediction appreciably.

\subsection*{Phase data}

\begin{table}[H]
    \centering
    \caption{Cursor readings and computed phase angles}
    \small
    \begin{tabular}{lrccr}
        \toprule
        Circuit & $f$ (Hz) & $T$ & $\Delta t$ & $\phi$ (\si{\degree}) \\
        \midrule
        RC, $V_C$   & 100   & \SI{10}{\milli\second}  & \SI{900}{\micro\second} & $-32.40$ \\
                    & 500   & \SI{2}{\milli\second}   & \SI{360}{\micro\second} & $-64.80$ \\
                    & 1000  & \SI{1}{\milli\second}   & \SI{240}{\micro\second} & $-86.40$ \\
                    & 5000  & \SI{200}{\micro\second} & \SI{48}{\micro\second}  & $-86.40$ \\
                    & 10000 & \SI{100}{\micro\second} & \SI{24}{\micro\second}  & $-86.40$ \\
        \midrule
        RL, $V_L$   & 100   & \SI{10}{\milli\second}  & \SI{480}{\micro\second} & $17.28$ \\
                    & 500   & \SI{2}{\milli\second}   & \SI{160}{\micro\second} & $28.80$ \\
                    & 1000  & \SI{1}{\milli\second}   & \SI{76}{\micro\second}  & $27.36$ \\
                    & 5000  & \SI{200}{\micro\second} & \SI{10}{\micro\second}  & $18.00$ \\
                    & 10000 & \SI{100}{\micro\second} & \SI{800}{\nano\second}  & $2.88$ \\
        \midrule
        RLC, $V_C$  & 100   & \SI{10}{\milli\second}  & \SI{520}{\micro\second} & $-18.72$ \\
                    & 500   & \SI{2}{\milli\second}   & \SI{300}{\micro\second} & $-54.00$ \\
                    & 1000  & \SI{1}{\milli\second}   & \SI{200}{\micro\second} & $-72.00$ \\
                    & 5000  & \SI{200}{\micro\second} & \SI{144}{\micro\second} & $-100.80$\,\textsuperscript{a} \\
                    & 10000 & \SI{100}{\micro\second} & \SI{68}{\micro\second}  & $-115.20$\,\textsuperscript{a} \\
        \midrule
        RLC, $V_L$  & 100   & \SI{10}{\milli\second}  & \SI{480}{\micro\second} & $17.28$ \\
                    & 500   & \SI{2}{\milli\second}   & \SI{320}{\micro\second} & $57.60$ \\
                    & 1000  & \SI{1}{\milli\second}   & \SI{212}{\micro\second} & $76.32$ \\
                    & 5000  & \SI{200}{\micro\second} & \SI{37}{\micro\second}  & $66.60$ \\
                    & 10000 & \SI{100}{\micro\second} & \SI{17}{\micro\second}  & $61.20$ \\
        \bottomrule
    \end{tabular}

    \vspace{0.4em}
    \footnotesize
    \raggedright
    \textsuperscript{a}Above resonance $\phi_C$ passes $-90^\circ$, and the cursor
    was placed on the following source crossing rather than the preceding one.
    The raw ratios give $259.20^\circ$ and $244.80^\circ$; subtracting
    $360^\circ$ recovers the values shown.
\end{table}

\subsection*{Series RC}

\begin{table}[H]
    \centering
    \caption{Measured values --- series RC ($R = \SI{4.7}{\kilo\ohm}$, $C = \SI{0.1}{\micro\farad}$)}
    \begin{tabular}{rccccc}
        \toprule
        & \multicolumn{2}{c}{$V_C$ (Vpp)} & \multicolumn{3}{c}{$\phi_C$ (\si{\degree})} \\
        \cmidrule(lr){2-3} \cmidrule(lr){4-6}
        $f$ (Hz) & Calculated & Measured & Calculated & Measured & Difference \\
        \midrule
        100   & 9.591 & 9.600 & $-16.45$ & $-32.40$ & $-15.95$ \\
        500   & 5.608 & 5.600 & $-55.89$ & $-64.80$ & $-8.91$  \\
        1000  & 3.207 & 3.600 & $-71.29$ & $-86.40$ & $-15.11$ \\
        5000  & 0.676 & 1.400 & $-86.13$ & $-86.40$ & $-0.27$  \\
        10000 & 0.338 & 0.064\,\textsuperscript{b} & $-88.06$ & $-86.40$ & $1.66$ \\
        \bottomrule
    \end{tabular}

    \vspace{0.4em}
    \footnotesize
    \raggedright
    \textsuperscript{b}Recorded value; the monotonic trend of the other four
    points suggests \SI{0.64}{\volt} peak-to-peak was intended.
\end{table}

The two lowest frequencies agree with the calculation to better than
\SI{0.2}{\percent} in amplitude. The three highest phase readings are identical
at $-86.40^\circ$, corresponding to $\Delta t / T = 0.24$ in all three cases,
which is correct at \SI{5}{\kilo\hertz} and \SI{10}{\kilo\hertz} but not at
\SI{1}{\kilo\hertz}, where the calculation calls for $\Delta t = \SI{198}{\micro\second}$.

\subsection*{Series RL}

The source amplitude at node A was recorded separately at each frequency because
the generator could not hold \SI{10}{\volt} peak-to-peak into this circuit. The
divider model uses $R_g = \SI{600}{\ohm}$ and an open-circuit amplitude of
\SI{10}{\volt} peak-to-peak.

\begin{table}[H]
    \centering
    \caption{Source amplitude at the circuit node --- series RL}
    \begin{tabular}{rccc}
        \toprule
        $f$ (Hz) & $|Z_{\text{circuit}}|$ (\si{\ohm}) & $V_s$ predicted (Vpp) & $V_s$ measured (Vpp) \\
        \midrule
        100   & 144.5  & 1.94 & 2.08 \\
        500   & 176.6  & 2.35 & 2.50 \\
        1000  & 251.9  & 3.27 & 3.40 \\
        5000  & 1046.5 & 8.21 & 8.20 \\
        10000 & 2078.1 & 9.44 & 9.40 \\
        \bottomrule
    \end{tabular}
\end{table}

Because the source amplitude varies with frequency, the amplitude comparison is
made on the ratio $V_L / V_s$, which is independent of $R_g$.

\begin{table}[H]
    \centering
    \caption{Measured values --- series RL ($R = \SI{100}{\ohm}$, $L = \SI{33}{\milli\henry}$)}
    \small
    \begin{tabular}{rcccccccc}
        \toprule
        & \multicolumn{2}{c}{measured (Vpp)} & \multicolumn{3}{c}{$V_L / V_s$}
        & \multicolumn{3}{c}{$\phi_L$ (\si{\degree})} \\
        \cmidrule(lr){2-3} \cmidrule(lr){4-6} \cmidrule(lr){7-9}
        $f$ (Hz) & $V_s$ & $V_L$ & Ideal & DCR & Meas. & Ideal & DCR & Meas. \\
        \midrule
        100   & 2.08 & 0.800 & 0.203 & 0.330 & 0.385 & $78.29$ & $17.49$ & $17.28$ \\
        500   & 2.50 & 1.600 & 0.720 & 0.635 & 0.640 & $43.97$ & $31.53$ & $28.80$ \\
        1000  & 3.40 & 2.030\,\textsuperscript{c} & 0.901 & 0.841 & 0.597 & $25.75$ & $22.88$ & $27.36$ \\
        5000  & 8.20 & 7.800 & 0.995 & 0.991 & 0.951 & $5.51$  & $5.48$  & $18.00$ \\
        10000 & 9.40 & 9.300 & 0.999 & 0.998 & 0.989 & $2.76$  & $2.76$  & $2.88$  \\
        \bottomrule
    \end{tabular}

    \vspace{0.4em}
    \footnotesize
    \raggedright
    \textsuperscript{c}This ratio is lower than the one recorded at
    \SI{500}{\hertz}. For a series RL circuit $V_L / V_s$ increases monotonically
    with frequency, so one of the two entries in this row is in error.
\end{table}

The ideal column and the measured column disagree by as much as $61^\circ$ at
\SI{100}{\hertz}, while the column that includes \SI{43}{\ohm} of winding
resistance agrees to $0.2^\circ$. The same holds for the amplitude ratio at
\SI{500}{\hertz} (0.635 predicted against 0.640 measured) and for the phase at
\SI{10}{\kilo\hertz}. The \SI{5}{\kilo\hertz} phase point is the exception: a
\SI{10}{\micro\second} cursor reading against a \SI{200}{\micro\second} period is
too coarse to resolve an angle whose true value is near \SI{3}{\micro\second}.

\subsection*{Series RLC}

\begin{table}[H]
    \centering
    \caption{Measured values --- series RLC ($R = \SI{1}{\kilo\ohm}$, $C = \SI{0.33}{\micro\farad}$, $L = \SI{10}{\milli\henry}$)}
    \small
    \begin{tabular}{rccccccccc}
        \toprule
        & \multicolumn{2}{c}{$V_C$ (Vpp)} & \multicolumn{2}{c}{$\phi_C$ (\si{\degree})}
        & \multicolumn{3}{c}{$V_L$ (Vpp)} & \multicolumn{2}{c}{$\phi_L$ (\si{\degree})} \\
        \cmidrule(lr){2-3} \cmidrule(lr){4-5} \cmidrule(lr){6-8} \cmidrule(lr){9-10}
        $f$ (Hz) & Calc. & Meas. & Calc. & Meas. & Ideal & DCR & Meas. & DCR & Meas. \\
        \midrule
        100   & 9.804 & 9.800 & $-11.73$  & $-18.72$  & 0.013 & 0.088 & 0.284 & $86.09$ & $17.28$ \\
        500   & 7.052 & 7.000 & $-46.98$  & $-54.00$  & 0.230 & 0.381 & 0.388 & $77.97$ & $57.60$ \\
        1000  & 4.447 & 4.400 & $-67.24$  & $-72.00$  & 0.579 & 0.677 & 0.680 & $77.52$ & $76.32$ \\
        5000  & 0.942 & 0.900 & $-102.28$ & $-100.80$ & 3.070 & 2.976 & 2.800 & $70.42$ & $66.60$ \\
        10000 & 0.417 & 0.440 & $-120.12$ & $-115.20$ & 5.435 & 5.277 & 4.860 & $57.00$ & $61.20$ \\
        \bottomrule
    \end{tabular}
\end{table}

The capacitor amplitudes match the calculation at all five frequencies, the
largest deviation being \SI{5.5}{\percent} at \SI{10}{\kilo\hertz}. The capacitor
phase readings above resonance confirm the sign reversal of $\angle Z_T$: the
measured $-100.80^\circ$ at \SI{5}{\kilo\hertz} and $-115.20^\circ$ at
\SI{10}{\kilo\hertz} sit within \SI{5}{\degree} of the calculated values, and both
lie beyond $-90^\circ$, which is only possible once the circuit has become
inductive. The inductor amplitudes again track the DCR column rather than the
ideal one, most clearly at \SI{1}{\kilo\hertz} where the measurement reads
\SI{0.680}{\volt} against \SI{0.677}{\volt} predicted with winding resistance and
\SI{0.579}{\volt} predicted without it. The inductor phase readings at
\SI{100}{\hertz} and \SI{500}{\hertz} do not reconcile with either model; the
\SI{100}{\hertz} entry duplicates the value recorded for the RL circuit at the
same frequency and is likely a transcription error.

%------------------------------------------------------------------------------
%------------------------------------------------------------------------------
\section*{Observations \& Discussion}

For the most part our measured values matched the theoretical calculations, but
the agreement was not uniform across the three circuits and the places where it
broke down turned out to be more instructive than the places where it held. The
RC circuit tracked theory the closest. The capacitor voltage at \SI{100}{\hertz}
and \SI{500}{\hertz} came out within a couple tenths of a percent of the
calculated amplitude, and the phase at the two lowest frequencies was within
about $9^\circ$. The one problem in the RC data is that the three highest phase
readings all came out to exactly $-86.40^\circ$, which corresponds to
$\Delta t / T = 0.24$ in every case. That value is right at \SI{5}{\kilo\hertz}
and \SI{10}{\kilo\hertz}, but at \SI{1}{\kilo\hertz} the calculation calls for
$\Delta t = \SI{198}{\micro\second}$ and a phase near $-71^\circ$, so getting
$-86^\circ$ there means the cursor most likely was not moved between those
measurements. It is a reading error rather than anything the circuit did.

The RL circuit is where the ideal model and the measurements disagreed the most,
and both major sources of error in this lab showed up in it. The first is the
function generator. The PB-503A does not hold \SI{10}{\volt} peak-to-peak into a
low-impedance load; at \SI{100}{\hertz} the source at the circuit node had sagged
to about \SI{2}{\volt} peak-to-peak. Fitting the measured source amplitude
against a voltage-divider model across all five RL frequencies gives an output
impedance of roughly \SI{600}{\ohm}, not the \SI{50}{\ohm} a bench generator is
usually assumed to have, with an rms residual of only \SI{0.07}{\volt}. Because
we probed channel~1 at the circuit node rather than at the generator terminals,
this \SI{600}{\ohm} sits outside the loop we measured and does not affect the
voltage ratios or the phase, which is why we compared $V_L / V_s$ rather than
$V_L$ alone. The second source is the inductor itself. An ideal \SI{33}{\milli\henry}
inductor has $|Z_L| = \SI{20.7}{\ohm}$ at \SI{100}{\hertz}, but a real inductor
also has winding resistance, and ours measured on the order of \SI{43}{\ohm}. At
low frequency that resistance is larger than the inductive reactance, so it
dominates the element voltage. The effect on the phase is dramatic: the ideal
calculation predicts $\phi_L = 78.3^\circ$ at \SI{100}{\hertz}, our measurement
gave $17.3^\circ$, and a model that includes the \SI{43}{\ohm} of series
resistance predicts $17.5^\circ$. The measured value sits within $0.2^\circ$ of
the corrected model, so the $61^\circ$ discrepancy is not error at all; it is the
winding resistance, which the ideal formula simply leaves out. The same
correction brings the amplitude ratios into line at the low-frequency end.

As we varied the frequency, the measured voltages moved exactly the way the
impedances predict. In the RC circuit the capacitor voltage started high and fell
as frequency rose, because $|Z_C| = 1/\omega C$ shrinks while the resistor stays
fixed, so a smaller and smaller share of the source appears across the capacitor.
This is low-pass behavior, and the phase confirmed it, moving from near $0^\circ$
at \SI{100}{\hertz} toward $-90^\circ$ at the high end as the circuit became
almost purely resistive from the capacitor's point of view. The RL circuit did
the opposite. The inductor voltage started small and climbed toward the full
source voltage as frequency rose, because $|Z_L| = \omega L$ grows with
frequency, so more of the source drops across the inductor. That is high-pass
behavior, and its phase moved from a large positive angle toward $0^\circ$ as the
inductor came to dominate the loop. The two reactive elements are mirror images
of each other, which is the whole point of running the same sweep on both.

The RLC circuit was the most interesting because it does something neither of the
single-reactance circuits can. Its resonant frequency is
$f_0 = 1/(2\pi\sqrt{LC}) = \SI{2.77}{\kilo\hertz}$, which falls between our
\SI{1}{\kilo\hertz} and \SI{5}{\kilo\hertz} test points, so the total reactance
changes sign as the sweep passes through it: the circuit is capacitive below
$f_0$ and inductive above it. Our measured capacitor phase caught this crossover.
Below resonance $\phi_C$ stayed between $0^\circ$ and $-90^\circ$, but at
\SI{5}{\kilo\hertz} and \SI{10}{\kilo\hertz} it read $-100.8^\circ$ and
$-115.2^\circ$, both beyond $-90^\circ$, which can only happen once the net
reactance has gone inductive and pushed the current phase past the point a purely
capacitive divider could reach. Those two points also illustrate the wrap problem
with the cursor method: the raw ratios gave $259.2^\circ$ and $244.8^\circ$
because we placed the cursor on the source crossing after the element crossing
rather than before it, and subtracting $360^\circ$ recovers the correct values.
The inductor amplitudes in this circuit again followed the model that includes
winding resistance rather than the ideal one, most clearly at \SI{1}{\kilo\hertz}
where we measured \SI{0.680}{\volt} against \SI{0.677}{\volt} predicted with
resistance and \SI{0.579}{\volt} without it. The low-frequency inductor phase
readings in the RLC circuit did not reconcile with either model, and the
\SI{100}{\hertz} entry is identical to the RL circuit's \SI{100}{\hertz} value,
which points to a transcription slip in the notebook rather than a real
measurement.

The output waveforms themselves did not change shape as we swept the frequency.
They stayed sinusoidal at the same frequency as the source throughout, which is
what we expect from linear R, L, and C elements: a sinusoid driven into a linear
network comes out a sinusoid at the same frequency, changed only in amplitude and
phase. What varied was the amplitude of the element voltage and its phase
relative to the source, and, in the RL circuit, the amplitude of the source
itself once the generator's output impedance came into play. We saw no clipping
or distortion, which tells us the generator was never driven into overload even
when the RL load pulled its output down to a couple of volts.

Overall the results confirmed the main ideas of the lab. Impedance behaves like
resistance in that the series and divider rules carry straight over, but unlike
resistance it is complex and frequency-dependent, so the element voltages come out
shifted in phase and the divide between them shifts as the frequency changes.
Where our numbers diverged from the ideal calculation, the difference was almost
always one of two real effects the ideal model leaves out --- the generator's
\SI{600}{\ohm} output impedance and the inductor's \SI{43}{\ohm} winding
resistance --- and once those are folded in, the measurements line up with the
theory.

%==============================================================================
\end{document}